
\documentclass{article}
\usepackage{colortbl}
\usepackage{makecell}
\usepackage{multirow}
\usepackage{supertabular}

\begin{document}

\newcounter{utterance}

\centering \large Interaction Transcript for game `dond', experiment `semi', episode 4 with Qwen3.6{-}35B{-}A3B{-}FP8{-}without{-}reasoning.
\vspace{24pt}

{ \footnotesize  \setcounter{utterance}{1}
\setlength{\tabcolsep}{0pt}
\begin{supertabular}{c@{$\;$}|p{.15\linewidth}@{}p{.15\linewidth}p{.15\linewidth}p{.15\linewidth}p{.15\linewidth}p{.15\linewidth}}
    \# & $\;$A & \multicolumn{4}{c}{Game Master} & $\;\:$B\\
    \hline

    \theutterance \stepcounter{utterance}  
    & & \multicolumn{4}{p{0.6\linewidth}}{
        \cellcolor[rgb]{0.9,0.9,0.9}{
            \makecell[{{p{\linewidth}}}]{
                \texttt{\tiny{[P1$\langle$GM]}}
                \texttt{Pelaat yhteistoiminnallista neuvottelupeliä, jossa sinun on sovittava toisen pelaajan kanssa siitä, kuinka jakaa joukko kohteita teidän välillä.} \\
\\ 
\texttt{Säännöt:} \\
\texttt{(a) Sinä ja toinen pelaaja saatte joukon kohteita. Kummallekin annetaan myös salainen arvofunktio, joka kuvaa, kuinka paljon arvostit kunkin kohteen tyyppiä.} \\
\texttt{(b) Vaihdat viestejä toisen pelaajan kanssa sopimaan siitä, kuka saa mitkä kohteet. Jokainen voi lähettää enintään 5 viestiä tai lopettaa pelin aikaisin tekemällä salaisen ehdotuksen milloin tahansa.} \\
\texttt{(c) Sinua pyydetään lähettämään salainen ehdotus, joka osoittaa kohteet, jotka haluat, muotoiltuna hakasulkeiden sisällä seuraavasti: "{[}Ehdotus: <numero> <kohteen nimi>, <numero> <kohteen nimi>, <...>{]}"} \\
\texttt{(d) Jos ehdotuksesi ovat täydentäviä, eli molempien ehdotusten täyttämiseksi on tarpeeksi kohteita, jokainen pelaaja saa pistemäärän saamiensa kohteiden arvojen summan perusteella. Muussa tapauksessa molemmat saavat nolla pistettä.} \\
\texttt{(e) Tavoitteesi on maksimoida pistemääräsi ja toisen pelaajan pistemäärän summa.} \\
\\ 
\texttt{Aloitetaan.} \\
\\ 
\texttt{Saatavilla olevien kohteiden joukko on:} \\
\\ 
\texttt{3 kirjat, 1 laukku, 1 pultti} \\
\\ 
\texttt{Salainen arvofunktiosi on:} \\
\\ 
\texttt{kirja: 0, laukku: 2, pultti: 8} \\
\\ 
\texttt{TÄRKEÄÄ: Viestit, paitsi salaiset ehdotukset, lähetetään suoraan toiselle pelaajalle, joten älä sisällytä mitään vastausta sääntöihin tai tekstiä, joka ilmoittaa viestistäsi. Salaisesti ehdotuksen tekemiseksi käytä ilmoitettua muotoa. Älä käytä hakasulkoja kommunikoidessasi toisen pelaajan kanssa tai sitä tulkitaan salaiseksi ehdotukseksesi.} \\
            }
        }
    }
    & \\ \\

    \theutterance \stepcounter{utterance}  
    & \multicolumn{4}{p{0.6\linewidth}}{
        \cellcolor[rgb]{0.9,0.9,0.9}{
            \makecell[{{p{\linewidth}}}]{
                \texttt{\tiny{[P1$\rangle$GM]}}
                \texttt{Hei!} \\
\\ 
\texttt{Ymmärsin pelin säännöt ja arvoni kohteille.} \\
\\ 
\texttt{Arvoni:} \\
\texttt{{-} Kirja: 0} \\
\texttt{{-} Laukku: 2} \\
\texttt{{-} Pultti: 8} \\
\\ 
\texttt{Koska arvostan pulttia eniten (8 pistettä) ja laukkua kohtalaisesti (2 pistettä), kun taas kirjat eivät tuo minulle arvoa (0 pistettä), ehdotan, että jaamme kohteet siten, että minä saan pultin ja laukun, ja sinä saat kaikki kirjat. Tämä maksimoi yhteisen arvon.} \\
\\ 
\texttt{Oletko valmis tekemään salaisen ehdotuksesi?} \\
            }
        }
    }
    & & \\ \\

    \theutterance \stepcounter{utterance}  
    & & & \multicolumn{2}{p{0.3\linewidth}}{
        \cellcolor[rgb]{0.9,0.9,0.9}{
            \makecell[{{p{\linewidth}}}]{
                \texttt{\tiny{[GM$|$GM]}}
                \texttt{continue} \\
            }
        }
    }
    & & \\ \\

    \theutterance \stepcounter{utterance}  
    & & & \multicolumn{2}{p{0.3\linewidth}}{
        \cellcolor[rgb]{0.9,0.9,0.9}{
            \makecell[{{p{\linewidth}}}]{
                \texttt{\tiny{[GM$|$GM]}}
                \texttt{Hei!} \\
\\ 
\texttt{Ymmärsin pelin säännöt ja arvoni kohteille.} \\
\\ 
\texttt{Arvoni:} \\
\texttt{{-} Kirja: 0} \\
\texttt{{-} Laukku: 2} \\
\texttt{{-} Pultti: 8} \\
\\ 
\texttt{Koska arvostan pulttia eniten (8 pistettä) ja laukkua kohtalaisesti (2 pistettä), kun taas kirjat eivät tuo minulle arvoa (0 pistettä), ehdotan, että jaamme kohteet siten, että minä saan pultin ja laukun, ja sinä saat kaikki kirjat. Tämä maksimoi yhteisen arvon.} \\
\\ 
\texttt{Oletko valmis tekemään salaisen ehdotuksesi?} \\
            }
        }
    }
    & & \\ \\

    \theutterance \stepcounter{utterance}  
    & & \multicolumn{4}{p{0.6\linewidth}}{
        \cellcolor[rgb]{0.9,0.9,0.9}{
            \makecell[{{p{\linewidth}}}]{
                \texttt{\tiny{[GM$\rangle$P2]}}
                \texttt{Pelaat yhteistoiminnallista neuvottelupeliä, jossa sinun on sovittava toisen pelaajan kanssa siitä, kuinka jakaa joukko kohteita teidän välillä.} \\
\\ 
\texttt{Säännöt:} \\
\texttt{(a) Sinä ja toinen pelaaja saatte joukon kohteita. Kummallekin annetaan myös salainen arvofunktio, joka kuvaa, kuinka paljon arvostit kunkin kohteen tyyppiä.} \\
\texttt{(b) Vaihdat viestejä toisen pelaajan kanssa sopimaan siitä, kuka saa mitkä kohteet. Jokainen voi lähettää enintään 5 viestiä tai lopettaa pelin aikaisin tekemällä salaisen ehdotuksen milloin tahansa.} \\
\texttt{(c) Sinua pyydetään lähettämään salainen ehdotus, joka osoittaa kohteet, jotka haluat, muotoiltuna hakasulkeiden sisällä seuraavasti: "{[}Ehdotus: <numero> <kohteen nimi>, <numero> <kohteen nimi>, <...>{]}"} \\
\texttt{(d) Jos ehdotuksesi ovat täydentäviä, eli molempien ehdotusten täyttämiseksi on tarpeeksi kohteita, jokainen pelaaja saa pistemäärän saamiensa kohteiden arvojen summan perusteella. Muussa tapauksessa molemmat saavat nolla pistettä.} \\
\texttt{(e) Tavoitteesi on maksimoida pistemääräsi ja toisen pelaajan pistemäärän summa.} \\
\\ 
\texttt{Aloitetaan.} \\
\\ 
\texttt{Saatavilla olevien kohteiden joukko on:} \\
\\ 
\texttt{3 kirjat, 1 laukku, 1 pultti} \\
\\ 
\texttt{Salainen arvofunktiosi on:} \\
\\ 
\texttt{kirja: 1, laukku: 2, pultti: 5} \\
\\ 
\texttt{TÄRKEÄÄ: Viestit, paitsi salaiset ehdotukset, lähetetään suoraan toiselle pelaajalle, joten älä sisällytä mitään vastausta sääntöihin tai tekstiä, joka ilmoittaa viestistäsi. Salaisesti ehdotuksen tekemiseksi käytä ilmoitettua muotoa. Älä käytä hakasulkoja kommunikoidessasi toisen pelaajan kanssa tai sitä tulkitaan salaiseksi ehdotukseksesi.} \\
\\ 
\texttt{Hei!} \\
\\ 
\texttt{Ymmärsin pelin säännöt ja arvoni kohteille.} \\
\\ 
\texttt{Arvoni:} \\
\texttt{{-} Kirja: 0} \\
\texttt{{-} Laukku: 2} \\
\texttt{{-} Pultti: 8} \\
\\ 
\texttt{Koska arvostan pulttia eniten (8 pistettä) ja laukkua kohtalaisesti (2 pistettä), kun taas kirjat eivät tuo minulle arvoa (0 pistettä), ehdotan, että jaamme kohteet siten, että minä saan pultin ja laukun, ja sinä saat kaikki kirjat. Tämä maksimoi yhteisen arvon.} \\
\\ 
\texttt{Oletko valmis tekemään salaisen ehdotuksesi?} \\
            }
        }
    }
    & \\ \\

    \theutterance \stepcounter{utterance}  
    & & & \multicolumn{4}{p{0.6\linewidth}}{
        \cellcolor[rgb]{0.9,0.9,0.9}{
            \makecell[{{p{\linewidth}}}]{
                \texttt{\tiny{[GM$\langle$P2]}}
                \texttt{Kiitos ehdotuksestasi! Huomasin kuitenkin, että mainitsit arvosi hieman eri tavalla kuin alkuperäisessä kuvauksessa. Kerroit arvosi olevan kirja: 0, laukku: 2, pultti: 8.} \\
\\ 
\texttt{Minun arvoni on kuitenkin:} \\
\texttt{{-} Kirja: 1} \\
\texttt{{-} Laukku: 2} \\
\texttt{{-} Pultti: 5} \\
\\ 
\texttt{Tämä tarkoittaa, että minulle kirjat ovat arvokkaita (1 kpl), laukku on keskitasoa (2) ja pultti on arvokkain (5).} \\
\\ 
\texttt{Yhteisen arvon maksimoimiseksi kannattaa jakaa kohteet niin, että kukin saa ne, joita hän arvostaa eniten.} \\
\\ 
\texttt{{-} Sinulle pultti on arvokkain (8), joten sinun kannattaa saada se.} \\
\texttt{{-} Minulle kirjat ovat arvokkaita (1 kpl x 3 = 3 pistettä), joten minun kannattaa saada kirjat.} \\
\texttt{{-} Laukku on arvokas molemmille (sinulle 2, minulle 2). Koska arvo on sama, voimme jakaa sen kumpaan tahansa, mutta jotta saamme molemmat tyytyväiseksi, voimme antaa sen sinulle, koska sinulla on jo pultti, tai minulle.} \\
\\ 
\texttt{Tarkastellaan kahta vaihtoehtoa:} \\
\\ 
\texttt{**Vaihtoehto 1:**} \\
\texttt{{-} Sinä: pultti (8) + laukku (2) = 10 pistettä} \\
\texttt{{-} Minä: 3 kirjaa (1+1+1) = 3 pistettä} \\
\texttt{{-} Yhteensä: 13 pistettä} \\
\\ 
\texttt{**Vaihtoehto 2:**} \\
\texttt{{-} Sinä: pultti (8) = 8 pistettä} \\
\texttt{{-} Minä: 3 kirjaa (3) + laukku (2) = 5 pistettä} \\
\texttt{{-} Yhteensä: 13 pistettä} \\
\\ 
\texttt{Molemmat vaihtoehdot tuottavat saman yhteisarvon (13 pistettä). Kuitenkin, jos haluat maksimoida oman pistemääräsi, vaihtoehto 1 on sinulle parempi (10 vs 8} \\
            }
        }
    }
     \\ \\

    \theutterance \stepcounter{utterance}  
    & & & \multicolumn{2}{p{0.3\linewidth}}{
        \cellcolor[rgb]{0.9,0.9,0.9}{
            \makecell[{{p{\linewidth}}}]{
                \texttt{\tiny{[GM$|$GM]}}
                \texttt{continue} \\
            }
        }
    }
    & & \\ \\

    \theutterance \stepcounter{utterance}  
    & & & \multicolumn{2}{p{0.3\linewidth}}{
        \cellcolor[rgb]{0.9,0.9,0.9}{
            \makecell[{{p{\linewidth}}}]{
                \texttt{\tiny{[GM$|$GM]}}
                \texttt{Kiitos ehdotuksestasi! Huomasin kuitenkin, että mainitsit arvosi hieman eri tavalla kuin alkuperäisessä kuvauksessa. Kerroit arvosi olevan kirja: 0, laukku: 2, pultti: 8.} \\
\\ 
\texttt{Minun arvoni on kuitenkin:} \\
\texttt{{-} Kirja: 1} \\
\texttt{{-} Laukku: 2} \\
\texttt{{-} Pultti: 5} \\
\\ 
\texttt{Tämä tarkoittaa, että minulle kirjat ovat arvokkaita (1 kpl), laukku on keskitasoa (2) ja pultti on arvokkain (5).} \\
\\ 
\texttt{Yhteisen arvon maksimoimiseksi kannattaa jakaa kohteet niin, että kukin saa ne, joita hän arvostaa eniten.} \\
\\ 
\texttt{{-} Sinulle pultti on arvokkain (8), joten sinun kannattaa saada se.} \\
\texttt{{-} Minulle kirjat ovat arvokkaita (1 kpl x 3 = 3 pistettä), joten minun kannattaa saada kirjat.} \\
\texttt{{-} Laukku on arvokas molemmille (sinulle 2, minulle 2). Koska arvo on sama, voimme jakaa sen kumpaan tahansa, mutta jotta saamme molemmat tyytyväiseksi, voimme antaa sen sinulle, koska sinulla on jo pultti, tai minulle.} \\
\\ 
\texttt{Tarkastellaan kahta vaihtoehtoa:} \\
\\ 
\texttt{**Vaihtoehto 1:**} \\
\texttt{{-} Sinä: pultti (8) + laukku (2) = 10 pistettä} \\
\texttt{{-} Minä: 3 kirjaa (1+1+1) = 3 pistettä} \\
\texttt{{-} Yhteensä: 13 pistettä} \\
\\ 
\texttt{**Vaihtoehto 2:**} \\
\texttt{{-} Sinä: pultti (8) = 8 pistettä} \\
\texttt{{-} Minä: 3 kirjaa (3) + laukku (2) = 5 pistettä} \\
\texttt{{-} Yhteensä: 13 pistettä} \\
\\ 
\texttt{Molemmat vaihtoehdot tuottavat saman yhteisarvon (13 pistettä). Kuitenkin, jos haluat maksimoida oman pistemääräsi, vaihtoehto 1 on sinulle parempi (10 vs 8} \\
            }
        }
    }
    & & \\ \\

    \theutterance \stepcounter{utterance}  
    & & \multicolumn{4}{p{0.6\linewidth}}{
        \cellcolor[rgb]{0.9,0.9,0.9}{
            \makecell[{{p{\linewidth}}}]{
                \texttt{\tiny{[P1$\langle$GM]}}
                \texttt{Kiitos ehdotuksestasi! Huomasin kuitenkin, että mainitsit arvosi hieman eri tavalla kuin alkuperäisessä kuvauksessa. Kerroit arvosi olevan kirja: 0, laukku: 2, pultti: 8.} \\
\\ 
\texttt{Minun arvoni on kuitenkin:} \\
\texttt{{-} Kirja: 1} \\
\texttt{{-} Laukku: 2} \\
\texttt{{-} Pultti: 5} \\
\\ 
\texttt{Tämä tarkoittaa, että minulle kirjat ovat arvokkaita (1 kpl), laukku on keskitasoa (2) ja pultti on arvokkain (5).} \\
\\ 
\texttt{Yhteisen arvon maksimoimiseksi kannattaa jakaa kohteet niin, että kukin saa ne, joita hän arvostaa eniten.} \\
\\ 
\texttt{{-} Sinulle pultti on arvokkain (8), joten sinun kannattaa saada se.} \\
\texttt{{-} Minulle kirjat ovat arvokkaita (1 kpl x 3 = 3 pistettä), joten minun kannattaa saada kirjat.} \\
\texttt{{-} Laukku on arvokas molemmille (sinulle 2, minulle 2). Koska arvo on sama, voimme jakaa sen kumpaan tahansa, mutta jotta saamme molemmat tyytyväiseksi, voimme antaa sen sinulle, koska sinulla on jo pultti, tai minulle.} \\
\\ 
\texttt{Tarkastellaan kahta vaihtoehtoa:} \\
\\ 
\texttt{**Vaihtoehto 1:**} \\
\texttt{{-} Sinä: pultti (8) + laukku (2) = 10 pistettä} \\
\texttt{{-} Minä: 3 kirjaa (1+1+1) = 3 pistettä} \\
\texttt{{-} Yhteensä: 13 pistettä} \\
\\ 
\texttt{**Vaihtoehto 2:**} \\
\texttt{{-} Sinä: pultti (8) = 8 pistettä} \\
\texttt{{-} Minä: 3 kirjaa (3) + laukku (2) = 5 pistettä} \\
\texttt{{-} Yhteensä: 13 pistettä} \\
\\ 
\texttt{Molemmat vaihtoehdot tuottavat saman yhteisarvon (13 pistettä). Kuitenkin, jos haluat maksimoida oman pistemääräsi, vaihtoehto 1 on sinulle parempi (10 vs 8} \\
            }
        }
    }
    & \\ \\

    \theutterance \stepcounter{utterance}  
    & \multicolumn{4}{p{0.6\linewidth}}{
        \cellcolor[rgb]{0.9,0.9,0.9}{
            \makecell[{{p{\linewidth}}}]{
                \texttt{\tiny{[P1$\rangle$GM]}}
                \texttt{Kiitos selkeytyksestä ja arvostuksistasi!} \\
\\ 
\texttt{Ymmärsin nyt tilanteen paremmin. Koska tavoitteenamme on maksimoida **yhteinen arvo** (sinun ja minun pistemääräni summa), molemmat vaihtoehdot ovat yhtä hyviä yhteiselle tulokselle (13 pistettä).} \\
\\ 
\texttt{Kuitenkin, jos tarkastelemme reilua jakoa tai minimoidaan eron meidän välillämme, vaihtoehto 2 on tasapuolisempi (minä 8, sinä 5, ero 3) verrattuna vaihtoehtoon 1 (minä 10, sinä 3, ero 7). Mutta koska sääntöjen mukaan tavoite on maksimoida summa, kumpi tahansa toimii.} \\
\\ 
\texttt{Teemme nyt molemmat salaiset ehdotukset samanaikaisesti varmistaaksemme, että ne ovat täydentäviä.} \\
\\ 
\texttt{Tässä on minun salainen ehdotukseni:} \\
\\ 
\texttt{{[}Ehdotus: 1 pultti, 1 laukku{]}} \\
            }
        }
    }
    & & \\ \\

    \theutterance \stepcounter{utterance}  
    & & & \multicolumn{2}{p{0.3\linewidth}}{
        \cellcolor[rgb]{0.9,0.9,0.9}{
            \makecell[{{p{\linewidth}}}]{
                \texttt{\tiny{[GM$|$GM]}}
                \texttt{proposal} \\
            }
        }
    }
    & & \\ \\

    \theutterance \stepcounter{utterance}  
    & & & \multicolumn{2}{p{0.3\linewidth}}{
        \cellcolor[rgb]{0.9,0.9,0.9}{
            \makecell[{{p{\linewidth}}}]{
                \texttt{\tiny{[GM$|$GM]}}
                [0, 1, 1]
            }
        }
    }
    & & \\ \\

    \theutterance \stepcounter{utterance}  
    & & \multicolumn{4}{p{0.6\linewidth}}{
        \cellcolor[rgb]{0.9,0.9,0.9}{
            \makecell[{{p{\linewidth}}}]{
                \texttt{\tiny{[GM$\rangle$P2]}}
                \texttt{Toinen pelaaja on jättänyt salaisen ehdotuksen. Nyt on sinun vuorosi jättää oma ehdotuksesi.} \\
            }
        }
    }
    & \\ \\

    \theutterance \stepcounter{utterance}  
    & & & \multicolumn{4}{p{0.6\linewidth}}{
        \cellcolor[rgb]{0.9,0.9,0.9}{
            \makecell[{{p{\linewidth}}}]{
                \texttt{\tiny{[GM$\langle$P2]}}
                \texttt{{[}Ehdotus: 1 kirja, 2 kirja, 3 kirja, 1 laukku{]}} \\
            }
        }
    }
     \\ \\

    \theutterance \stepcounter{utterance}  
    & & & \multicolumn{2}{p{0.3\linewidth}}{
        \cellcolor[rgb]{0.9,0.9,0.9}{
            \makecell[{{p{\linewidth}}}]{
                \texttt{\tiny{[GM$|$GM]}}
                \texttt{proposal must include every item type only once} \\
            }
        }
    }
    & & \\ \\

\end{supertabular}
}

\end{document}
